Huntington-Hill is U.S.\ law.
Since November~15, 1941, 2~U.S.C.\ \S2a has required that congressional
seats be apportioned among the states using the method of equal proportions,
which is the formal name for Huntington-Hill.
Every redistricting cycle --- 1950, 1960, \ldots, 2020 --- has begun with
a Huntington-Hill apportionment.
The method's outputs determine how many congressional districts each state
draws: North Carolina drew 14 districts after the 2020 apportionment;
Montana, which received a second seat for the first time since 1990,
drew 2.

Despite its statutory status, Huntington-Hill is often poorly understood
by practitioners and even by redistricting attorneys.
The geometric mean rounding rule --- ``round at $\sqrt{s(s+1)}$'' ---
is less intuitive than Webster's arithmetic mean rounding (``round at
$s+0.5$'') or Hamilton's largest-remainder rule.
The priority formula $P_i(s) = p_i/\sqrt{s(s+1)}$ is sometimes
described informally as ``give the next seat to whichever state needs
it most,'' but this informal description does not convey the specific
geometric mean construction or its optimality property.

This paper develops the mathematical foundations of Huntington-Hill
with enough formal detail for use in expert witness testimony,
policy analysis, and software implementation verification.
Section~\ref{sec:definition} gives the precise mathematical definition.
Section~\ref{sec:properties} proves the two key properties:
optimality (HH minimises the population pair test) and paradox immunity.
Section~\ref{sec:comparison} compares HH to Webster and Adams,
the two closest neighbours on the divisor method spectrum.
Section~\ref{sec:history} documents the 1941 statutory adoption
and the congressional debate that preceded it.
Section~\ref{sec:montana} analyses \textit{United States Dept.\ of
Commerce v.\ Montana} (1992), the definitive Supreme Court ruling on
Huntington-Hill's constitutionality.
Section~\ref{sec:implementation} documents the current \textsc{bisect-apportion}
implementation boundary and the Census-reference fixtures still needed for
SHA-verified official-output claims.

\subsection{Data}

All empirical results use the 2020 Census apportionment populations
published by the Census Bureau on April~26, 2021
\citep{census2020apportionment}.
Total apportionment population: 331{,}108{,}434.
House size: 435.
The apportionment population excludes the District of Columbia
but includes U.S.\ service members and federal civilian employees
stationed abroad who are allocated to a home state.
